%% ch13_conclusions.tex
%% Chapter 13: Conclusions and Future Work
%% Thesis: Self-Adaptive Model-Based Protection for Future Power Systems
%% ============================================================
\chapter{Conclusions and Future Work}
\label{ch:conclusions}

\section{Summary of Contributions}
\label{sec:ch13_summary}

This thesis has developed, analysed, and validated the \textbf{Self-Adaptive
Model-Based Protection (SAMBP)} framework---a hierarchical, physics-aware
architecture for protection of future power systems dominated by
Inverter-Based Resources (IBRs). The framework treats adaptive protection
as an \emph{inverse estimation problem}: each protection element maintains
an explicit physical model of its protected device and continuously
identifies the current model parameters from real-time measurements.
Adaptation decisions are gated by a four-component confidence score and
bounded by a provably safe update law.

The thesis delivers 25 original contributions across three tiers, of
which 17 are fully validated with numerical results ([P]) and 8 represent
complete theoretical frameworks with placeholder results pending the
18--24 month experimental execution phase ([F]).

\subsection{Tier 1: Local Relay Physics}

\paragraph{Chapter~\ref{ch:line_diff} --- Line Differential Protection (C1--C5).}
Five adaptive 87L elements provide complementary coverage across all
fault types. The physics-constrained adaptive threshold $k_m^*(\kibr)$
(C1) achieves TPR = 1.000 across $\kibr \in [0.03, 0.85]$~pu. The
87LN negative-sequence element (C3) proves that IBR negative-sequence
suppression \emph{improves} 87LN sensitivity ($\rho_2 = 2$, independent
of $\kibr$). The 87LN0 zero-sequence element (C5) provides the only
viable SLG protection on pure-IBR feeders with no SG grounding.

\paragraph{Chapter~\ref{ch:generator_oc} --- Generator Overcurrent (C6--C10).}
The SAMBP-SyncOC framework introduces the first inverse-estimation-based
adaptive OC protection for synchronous generators. The two-pass LM
estimator (C6) achieves $R^2 > 0.999$ with Jacobian condition number
$\kappa_n(J) < 20$ across three-phase, line-to-line, and SLG faults.
The confidence-gated bounded update law (C7, C8) reduces spurious pickups
by $\approx 60\%$ with no security degradation. Extensions to HIF
detection, sequence estimation, multi-generator buses, Park $dq$ truth
model validation, and DFIG Type-3 wind generators (C9, C10) are
theoretically complete and pending experimental execution.

\paragraph{Chapter~\ref{ch:transformer_diff} --- Transformer Differential (C11--C12).}
The 87T-SPRT scheme replaces the ad-hoc second harmonic blocking criterion
with a Wald SPRT on the harmonic ratio trajectory, providing explicit
false-trip and miss-rate guarantees ($\alpha = \beta = 0.001$). The
IBR-adaptive bias slope $k_T^*(\kibr)$ closes the second failure mode
(false restrain during IBR-fed internal fault). Both contributions are
theoretically complete, with expected decision time $\approx 50$~ms
comparable to classical blocking delay. Experimental validation is pending.

\paragraph{Chapter~\ref{ch:distance} --- Distance Protection (C13--C14).}
The quadrilateral 21 characteristic with IBR reactive-shift correction (C13)
achieves $P_\mathrm{detect} = 99.3\%$, $P_\mathrm{secure} = 99.97\%$.
The Monte Carlo reliability framework (C14) provides the first probabilistic
reliability quantification for IBR distance protection.

\paragraph{Chapter~\ref{ch:microgrid} --- Microgrid Protection (C15--C17).}
The GFL/GFM dual-mode characterisation (C15) quantifies the 4--8$\times$
fault current reduction at grid-to-island transition. The ROCOF islanding
detection (C16) achieves $< 1\%$ NDZ and 114~ms isolation time compliant
with IEEE~1547-2018. The adaptive OC setting-group (C17) switches within
5~ms of GOOSE receipt.

\subsection{Tier 2: Digital Monitoring and Learning}

\paragraph{Chapter~\ref{ch:digital_twin} --- Digital Twin (C18--C19).}
The non-blocking DT architecture (C18) guarantees that EKF updates never
delay primary trip decisions through the shadow register pattern. The
EKF joint estimation (C19) achieves RMSE $\leq 0.024$~pu for $\kibr$,
$\leq 0.5\%$ for $\Zline$, and $\leq 1.1\%$ for fault location, all
bounded by the analytically derived Cramr-Rao Lower Bound.

\paragraph{Chapter~\ref{ch:ml} --- Physics-Informed ML (C20--C21).}
The PINN classifier (C20) achieves 96.1\% five-class accuracy with zero
physics-constraint violations, by embedding three KVL/KCL constraints
in the loss function. CUSUM-SGD online learning (C21) adapts to fleet
composition drift within 21 detection samples with $< 2$~pp catastrophic
forgetting.

\subsection{Tier 3: Centralised Coordination and Security}

\paragraph{Chapter~\ref{ch:centralised} --- WAPC (C22--C24).}
Wide-area PMU fault localisation (C22) achieves 94.5\% correct-line
identification at $\sigma_V = 0.01$~pu. The CTI-graded settings engine
(C23) produces 100\% CTI-compliant settings in $< 3$~s. The closed-loop
CUSUM$\to$settings$\to$GOOSE protocol (C24) completes the full adaptive
cycle in $\approx 31$~minutes with 8 of 9 steps automated.

\paragraph{Chapter~\ref{ch:cybersecurity} --- Cybersecurity (C17, C25).}
HMAC-SHA256 (C17) reduces GOOSE attack risk by 54.9\% at 0.72~ms overhead.
The adversarial PINN robustness framework (C25) is theoretically established
and pending experimental validation.

\section{Mapping to Research Objectives}
\label{sec:ch13_objectives}

\begin{table}[htbp]
\caption{Research objectives and their fulfilment status}
\label{tab:ch13_objectives}
\small
\begin{tabularx}{\textwidth}{@{}lp{2cm}p{4cm}p{3cm}l@{}}
\toprule
Obj. & Chapter & Contribution & Key metric & Status \\
\midrule
O1 & \ref{ch:line_diff} & C1--C5: Adaptive 87L
  & TPR = 1.000, all $\kibr$ & [P] \\
O2 & \ref{ch:generator_oc} & C6--C8: SAMBP-SyncOC
  & $R^2 > 0.999$; $-60\%$ pickups & [P], C9-C10 [F] \\
O3 & \ref{ch:transformer_diff} & C11--C12: 87T-SPRT
  & $\alpha = \beta = 0.001$; 50~ms & [F] \\
O4 & \ref{ch:distance} & C13--C14: Quadrilateral 21
  & $P_d = 99.3\%$, $P_s = 99.97\%$ & [P] \\
O5 & \ref{ch:microgrid} & C15--C17: Microgrid
  & 114~ms, NDZ $< 1\%$ & [P] \\
O6 & \ref{ch:digital_twin} & C18--C19: EKF DT
  & RMSE $\leq 0.024$~pu & [P] \\
O7 & \ref{ch:ml} & C20--C21: PINN + CUSUM
  & 96.1\%; $< 2$~pp forgetting & [P] \\
O8 & \ref{ch:centralised},\ref{ch:cybersecurity} & C22--C25: WAPC + security
  & 94.5\%; 54.9\% risk reduction & [P], C25 [F] \\
\bottomrule
\end{tabularx}
\end{table}

\section{Limitations}
\label{sec:ch13_limitations}

\begin{enumerate}
\item \textbf{Validated [P] scope:} All [P] contributions are validated
      by EMT simulation and software-in-the-loop. Physical hardware-in-the-loop
      on commercial IEDs (SEL-411L) is pending.

\item \textbf{Synthetic training data:} PINN training and EKF validation
      data are synthetic. Real field fault recordings from live IBR
      substations are needed for independent validation (Gap G7).

\item \textbf{Single-substation WAPC:} The WAPC architecture is designed
      for up to 34 buses within one substation. Multi-substation
      coordination requires a hierarchical extension (Gap G4).

\item \textbf{GFM model horizon:} GFM inverters are modelled as VSMs
      valid for $t \leq 150$~ms. AVR dynamics beyond this horizon are
      not captured (Gap G6).

\item \textbf{Pending [F] contributions:} Eight contributions
      (C9, C10, C11, C12, C25 and coordinated attack analysis) are
      theoretically complete but require experimental execution. The
      simulation scripts exist and the theoretical frameworks are validated
      against Milestone-1 results.
\end{enumerate}

\section{Future Work and Research Roadmap}
\label{sec:ch13_future}

\subsection{Phase 1 (0--6 months): Complete Existing Scripts}

\paragraph{F1.1 --- Execute SAMBP-SyncOC Milestones 2--5.}
Run \texttt{run\_milestone2.py}, \texttt{run\_multi\_gen.py}, and
\texttt{run\_park\_dq\_validation.py} (all existing, structurally validated).
Fill in the 12 result placeholders in Section~\ref{sec:ch5_extensions}
and publish Contribution C9 as a journal paper.

\textbf{Target:} $R^2 > 0.999$ for HIF at $R_F \leq 100\,\Omega$;
multi-generator CTI compliance $> 95\%$;
Park $dq$ truth model agreement $< 5\%$ RMS at $t \leq 250$~ms.

\paragraph{F1.2 --- Execute 87T-SPRT Simulation.}
Implement the transformer energisation and turn-to-turn fault simulator.
Fit $(\mu_0, \sigma_0, \mu_1, \sigma_1)$ from simulation data.
Execute 400-scenario study plan of Section~\ref{sec:ch6_validation}.
Publish Contributions C11 and C12.

\textbf{Target:} False trip rate $\leq 0.1\%$ on inrush;
miss rate $\leq 0.1\%$ on internal fault;
$> 30\%$ improvement over classical 2nd harmonic blocking at $\kibr < 0.15$.

\subsection{Phase 2 (6--12 months): Extend to DFIG and Adversarial PINN}

\paragraph{F2.1 --- DFIG Type-3 Adaptive OC (Contribution C9d).}
Implement the piecewise-ODE DFIG crowbar model.
Apply two-phase LM estimator across 30 crowbar activation scenarios.
Validate $R^2 > 0.99$ and $\Delta t_{cb} < 2$~ms.

\paragraph{F2.2 --- Adversarial PINN Robustness (Contribution C25).}
Implement FGSM/PGD attack generators.
Validate $> 80\%$ adversarial accuracy at $\varepsilon = 0.05$~pu.
Prove or disprove the physics-constraint robustness hypothesis
(Eq.~\eqref{eq:ch12_margin_bound}).

\subsection{Phase 3 (12--18 months): GFM Extended Dynamics and HIL}

\paragraph{F3.1 --- GEAC Stability Boundary (Gap G6).}
Extend the GFM VSM model to include AVR dynamics beyond $t = 150$~ms.
Determine the Grid-Edge Adaptive Characteristic (GEAC) stability
boundary as a function of $(M, D, L_v)$ and $\kibr$.

\paragraph{F3.2 --- Hardware-in-the-Loop Validation (Gap G1).}
Commission HIL test cell with SEL-411L relay and RTDS real-time simulator.
Execute 100-scenario regression suite against physical relay firmware.

\subsection{Phase 4 (18--24 months): Multi-Substation and Field Data}

\paragraph{F4.1 --- Multi-Substation WAPC Scaling (Gap G4).}
Extend the WAPC to a hierarchical architecture: local substation
coordinators reporting to a regional coordinator.
Solve the consensus problem under 10--50~ms propagation delay.

\paragraph{F4.2 --- Field Data Validation (Gap G7).}
Collaborate with a DSO operating an IBR-dominated network to collect
and label real fault recordings. Validate PINN classifier and EKF
on real-world data. Quantify the synthetic-to-real domain gap.

\subsection{Summary Roadmap}

\begin{table}[htbp]
\caption{18--24 month research roadmap}
\label{tab:ch13_roadmap}
\small
\begin{tabularx}{\textwidth}{@{}lp{3cm}lXl@{}}
\toprule
Phase & Task & Timeline & Target metric & Contribs. \\
\midrule
1.1 & SyncOC Milestones 2--5 & 0--3 mo & $R^2>0.999$, CTI $>95\%$ & C9, C10 \\
1.2 & 87T-SPRT simulation & 3--6 mo & FTR $<0.1\%$, MR $<0.1\%$ & C11, C12 \\
2.1 & DFIG Type-3 OC & 6--9 mo & $R^2>0.99$, $|\Delta t_{cb}|<2$~ms & C9d \\
2.2 & Adversarial PINN & 9--12 mo & $>80\%$ adv.\ acc.\ & C25 \\
3.1 & GEAC stability & 12--15 mo & Stability boundary mapped & G6 \\
3.2 & HIL validation & 15--18 mo & 95/100 scenarios pass & G1 \\
4.1 & Multi-substation WAPC & 18--21 mo & Consensus $<100$~ms & G4 \\
4.2 & Field data validation & 21--24 mo & PINN accuracy $>90\%$ real & G7 \\
\bottomrule
\end{tabularx}
\end{table}

\section{Concluding Remarks}
\label{sec:ch13_concluding}

The displacement of synchronous machines by IBRs is irreversible and
accelerating. The protection crisis it creates is not a temporary
transition problem but a permanent structural feature of all future
high-IBR power systems. Legacy relay algorithms, designed for passive,
high-current, physics-determined fault current sources, fail systematically
when the fault current is controlled, software-limited, and compositionally
variable.

The SAMBP framework proposed in this thesis addresses this crisis from
first principles: by treating adaptive protection as an inverse estimation
problem, it provides a unified mathematical structure applicable to every
protection element---line differential, generator overcurrent, transformer
differential, distance, and islanding---and to every source type.
The framework is the first to integrate model identification, real-time
state estimation, and centralised coordination into a closed-loop self-adaptive
system with provable safety guarantees.

The 17 validated contributions establish that the framework works: relay
algorithms adapting in real-time to measured IBR fleet composition achieve
full detection coverage across the operating envelope, with quantified
reliability and security bounds. The 8 pending contributions extend the
framework to the broader technology space---DFIG wind generators,
transformer protection, adversarial cyberattacks---that defines the full
scope of future power system protection needs.

The remaining open gaps (multi-substation scaling, HIL firmware validation,
and real field fault data) define a clear post-doctoral research agenda.
This thesis has equipped that agenda with the foundational methodology,
validated algorithms, documented performance bounds, and existing simulation
scripts needed to make rapid progress on the 18--24 month timeline.
